%% vert-punct.tex — 直排标点宽度调整的 clreq 断言用例（仅由 clreq_test.py 编译）
%% 每一段构造一种上下文，供 test/clreq_test.py 的 run_vertical_assertions 度量：
%%   ① 汉字之间的单个标点占满一字幅（不再无条件挤压）
%%   ② 相邻标点（夹注符号参与）合计 1.5 字幅
%%   ③ 直排冒号/分号/问号/叹号固定一字幅
%%   ④ 挤压方向（收回哪一侧的空白，字面就往哪边让）
%%   ⑤ 行首禁则：让点号恰好越出列末，看它是否被挤进本列
\documentclass{ltc-cn-vbook}
\setmainfont{TW-Kai}

\begin{document}

\begin{正文}
天地玄黄，宇宙洪荒。

金木水火土。「引号内容」接续文字。

子曰：学而时习之；不亦说乎？有朋自远方来！

他说：「好。」「好。」连续引号。

% P1 扩类：中国大陆式间隔号占半个字宽（clreq 附录 A 间隔号 GB 式），
% 连接号 ～ 与分隔号 ／ 全角形式占满一字幅。
名词·解释间隔号。

范围～连接号；甲／乙分隔号。

% 避头相位扫描：数字串长度自 28 至 38 逐段递增，扫过列末的各个相位。
% 本版式每列约 33 字，其中必有若干段的逗号恰好越出列末——禁则须把它挤进
% 本列，而不是让它成为下一列的首字。用固定长度做这个用例是脆的（容量随
% 版式与字体变），扫相位才与容量无关。
一二三四五六七八九十一二三四五六七八九十一二三四五六七八，避头甲

一二三四五六七八九十一二三四五六七八九十一二三四五六七八九，避头乙

一二三四五六七八九十一二三四五六七八九十一二三四五六七八九十，避头丙

一二三四五六七八九十一二三四五六七八九十一二三四五六七八九十一，避头丁

一二三四五六七八九十一二三四五六七八九十一二三四五六七八九十一二，避头戊

一二三四五六七八九十一二三四五六七八九十一二三四五六七八九十一二三，避头己

一二三四五六七八九十一二三四五六七八九十一二三四五六七八九十一二三四，避头庚

一二三四五六七八九十一二三四五六七八九十一二三四五六七八九十一二三四五，避头辛

一二三四五六七八九十一二三四五六七八九十一二三四五六七八九十一二三四五六，避头壬

一二三四五六七八九十一二三四五六七八九十一二三四五六七八九十一二三四五六七，避头癸

一二三四五六七八九十一二三四五六七八九十一二三四五六七八九十一二三四五六七八，避头子

% 列首开始夹注符号相位扫描：「不得位于行末（clreq 行尾禁则），恰好越出列末
% 时必被推到下一列列首。此时它的始端空白应整段收回，字面紧贴列首——量法是
% 比较「列首为引号」与「列首为汉字」两列第二字的位置，两者应相差半个字幅。
% 同样用扫相位而不是固定长度，与列容量无关。
一二三四五六七八九十一二三四五六七八九十一二三四五六七八九十一二「列首引甲」接续文字。

一二三四五六七八九十一二三四五六七八九十一二三四五六七八九十一二三「列首引乙」接续文字。

一二三四五六七八九十一二三四五六七八九十一二三四五六七八九十一二三四「列首引丙」接续文字。

一二三四五六七八九十一二三四五六七八九十一二三四五六七八九十一二三四五「列首引丁」接续文字。

% 标号位置：脚注标号组被 cm 缩放绘制，解析器须跟踪 cm 才能读到真实坐标
% （只认单位矩阵 Tm 的旧解析会把 ︻一︼ 读到页面左上角的变换前原点）。
标号位置丙\脚注{注文。}丁戊己庚

\输出脚注
\end{正文}

\end{document}
